% =====================================================================================================================
\section{Introduction}
\label{sec:intro}

Real-time arrival information shortens the time passengers wait at stops and makes that wait less uncertain
\parencite{watkins2011,brakewood2019}, but only if it is accurate: inaccurate predictions erode trust and can make
passengers miss the bus they were told to expect \parencite{gooze2013benefits,rahman2013users}. Most passengers
receive these predictions from the operator itself, through its own app and stop displays, or from third-party apps
that relay the operator's figures. How accurate the operator's predictions are, and whether a third party with only
the operator's open data can do better, is rarely measured on the same buses and at the same moments.

This paper makes that measurement for Bizkaibus, the interurban bus network of the province of Biscay (Spain). Since
2023 Bizkaibus shows its own real-time waiting times in an official app and on some stop displays. It also publishes,
as open data, the GPS position of every vehicle in service (a SIRI Vehicle Monitoring feed, \cite{siri}) and its
timetable (GTFS, \cite{gtfs}), but not the predictions themselves. A companion paper
\parencite{garrido2026metodo} and a technical report \parencite{garrido2026memoria} describe and evaluate a method
that predicts arrivals from those two open sources alone, used by the independent Bizkaibus+ app. That evaluation
compared the method with the published timetable, which is the obvious but weak baseline: after the first hour of a
trip, most Bizkaibus buses pass a stop more than two minutes before the published time.

Here the baseline is the operator's own live prediction. For one week in November 2026 we recorded, every minute, the
waiting times that the official service gave for 60 stops, together with the live feed every 20 seconds. We then
compared, for the same bus at the same stop and the same minute, the waiting time announced by the operator with the
arrival predicted by the method, and both with the moment the bus actually passed.

The contributions are:
\begin{enumerate}
  \item A paired protocol to compare a third-party predictor with an operator's real-time predictions when the
        operator publishes positions but not predictions: what is captured, how each announcement is matched to a
        vehicle, how the actual passage is reconstructed and how the comparison avoids giving either side an
        advantage.
  \item The comparison itself, over one week of data, by how far ahead the bus is, by type of day and by type of stop,
        with confidence intervals that respect the dependence between predictions of the same trip.
  \item The data, the capture program and the analysis, openly available, so that the comparison can be repeated
        on other days or other networks.
\end{enumerate}

Section~\ref{sec:related} reviews related work, Section~\ref{sec:data} describes the data,
Section~\ref{sec:design} the comparison, Section~\ref{sec:results} the results and Section~\ref{sec:discussion} their
interpretation and limits.

% =====================================================================================================================
\section{Related work}
\label{sec:related}

The accuracy of operator real-time information has been studied mostly through the operator's own logs.
\textcite{cats2016realtime} evaluated the bus arrival information system of Stockholm against automatic vehicle
location records and found that prediction errors grow with the prediction horizon and depend on the type of stop and
on the time of day. Studies of passenger benefits \parencite{watkins2011,mishalani2006,liu2020waiting} use the
operator's predictions as given and show that the value of real-time information depends on its accuracy
\parencite{gooze2013benefits}. On the other side, third-party predictors built only on public data have been shown
to be feasible: \textcite{garciamaurino2024improving} improve arrival predictions using only data available through
public transport APIs, and \textcite{wessel2017retrospective} reconstruct actual passage times from a public vehicle
location feed, the same kind of ground truth used here. The gap that this paper addresses is the direct comparison:
operator predictions and a third-party predictor on the same vehicles, stops and moments, against the same actual
passage, with the operator's predictions recorded as passengers saw them.

The evaluation follows the rules of out-of-sample forecasting evaluation \parencite{tashman2000outofsample}: the
method is calibrated only with data from before the week that is measured, so that nothing from the test days leaks
into it \parencite{kaufman2012leakage}.

% =====================================================================================================================
\section{Data}
\label{sec:data}

\subsection{What was captured}

A program running on GitHub Actions captured, from 06:00 to 23:00 Madrid time (the night feed repeats stale data):
\begin{itemize}
  \item \textbf{The live feed}, every 20 seconds: the public SIRI Vehicle Monitoring file of Bizkaibus, with the GPS
        position, the trip and the next stop of every vehicle in service. The file changes every 60 seconds, so two
        of every three copies repeat the previous one.
  \item \textbf{The official waiting times}, every minute, for each of 60 stops: the response of the web service
        behind the official Bizkaibus app (\texttt{GetPasoParadaMobile\_JSON}). For each line that serves the stop,
        the service gives the next two buses (here $e_1$ and $e_2$), each with its waiting time in whole minutes and
        its distance in metres. The stops are queried one after another, so each stop is queried about once a minute.
  \item \textbf{The timetable} of the day (GTFS), once a day.
\end{itemize}

The 60 stops are intermediate stops with many buses and lines, at least 15 minutes after the start of the trip (where
predictions are hardest and most useful), plus 28 \emph{priority} stops: stops pointed out by passengers and the Bilbao
stops where the actual passage deviates most from the published time. Each stop and the reason for its choice are
listed in the repository (\texttt{paradas.json}).

\subsection{Days}

The week runs from Sunday 8 to Saturday 14 November 2026 (Table~\ref{tab:days}). Wednesday 11 was a general strike:
Bizkaibus ran minimum services (about 30\%), with about 160 vehicles in the feed against 240 on a normal day. That day
is analysed separately and does not enter the results of the week. Three days are incomplete: on Monday 9 the morning
block was lost because the capture overran its time limit; on Thursday 12 the official service stopped answering over
plain HTTP at 17:12 (it answers over HTTPS, which the capture uses from then on), so that afternoon has no official
times; and on Saturday 14 the capture started at 09:41. The live feed has no gap longer than two minutes on any day.

\begin{table}[t]
\centering
\caption{Days of the capture (Madrid time). The feed and the official service were captured from about 06:00 to
23:00, except as noted.}
\label{tab:days}
\small
\begin{tabularx}{\linewidth}{llX}
\toprule
Date & Day & Coverage \\
\midrule
8 Nov 2026 & Sunday & Complete \\
9 Nov 2026 & Monday & 12:42--23:00 (morning block lost: the capture overran its time limit) \\
10 Nov 2026 & Tuesday & Complete \\
11 Nov 2026 & Wednesday & Complete, \emph{analysed separately}: general strike, minimum services (about 30\%) \\
12 Nov 2026 & Thursday & Feed complete; official times until 17:12 (the service stopped answering over HTTP) \\
13 Nov 2026 & Friday & Complete \\
14 Nov 2026 & Saturday & From 09:41; to be added in the next version \\
\bottomrule
\end{tabularx}
\end{table}


% =====================================================================================================================
\section{Comparison design}
\label{sec:design}

\subsection{The two predictions}

\paragraph{Official} For a query at time $t$ to stop $s$, an announcement of $m$ minutes for a bus of line $\ell$
gives the predicted arrival $\hat a_{\mathrm{off}} = t + 60\,m$.

\paragraph{Method} The method of \textcite{garrido2026metodo} is recomputed from the captured feed with its open
reference implementation, which reproduces the production engine exactly (it passes the 407 test vectors produced by
that engine). It is not taken from the live app, so that anyone can repeat it. Its learned tables (segment times, pace
factors and per-trip corrections) are those of Saturday 7 November, calibrated only with data from before the week,
and they are not updated during the week. For the query at time $t$, the method's prediction is the one it gave in the
last copy of the feed before $t$, provided that copy is at most 60 seconds old; it is therefore never more recent
than the official one. Its predicted arrival $\hat a_{\mathrm{met}}$ is compared in seconds and also rounded to whole
minutes from $t$, as a passenger would see it:
$\hat a_{\mathrm{met,min}} = t + 60\,\max\bigl(0, \operatorname{round}\bigl((\hat a_{\mathrm{met}} - t)/60\bigr)\bigr)$.
The main comparison uses the rounded value, so that neither side gains from the resolution of its figures.

\subsection{The actual passage}

The actual passage of a bus at a stop, $a$, is reconstructed from the feed as in the companion paper: it is the midpoint
between the last copy of the feed in which the bus had not yet passed the stop (its next stop is the stop itself or an
earlier one) and the first copy in which it had, without the bus changing trip in between. Because the feed changes
every 60 seconds, $a$ is known to within about $\pm30$ seconds. The same $a$ is used for both predictions, so this
uncertainty adds the same noise to both errors and does not favour either.

\subsection{Matching announcements to buses}

The official service does not say which vehicle it announces. For a query at $t$ to stop $s$, its $e_1$ for line $\ell$
is matched to the first bus of line $\ell$ that actually passes $s$ after $t$, and its $e_2$ to the second. This is the
bus a passenger waiting at $s$ would take. Lines listed twice in the same response (two routes of the same line) are
left out, as are announcements whose matched bus passes more than 30 minutes after $t$ or is not tracked by the feed.
Because buses of the same line can bunch, the results are repeated without the pairs in which another bus of the same
line passes within two minutes of the matched one; the conclusions do not change (Section~\ref{sec:results}).

\subsection{Metrics}

The error of each prediction is $\hat a - a$ (positive: the bus arrives earlier than announced). For each source we
give the median absolute error, the share of predictions within $\pm1$ and $\pm2$ minutes, and the median signed
error, by horizon $a - t$ (0--3, 3--6, 6--10, 10--20 and 20--30 minutes), by type of day and by type of stop.
Predictions of the same trip are not independent, so confidence intervals for the difference between the two sources
are obtained by bootstrap over whole trips \parencite{efron1993bootstrap}: each of 400 resamples draws trips with
replacement, with all their pairs.

% =====================================================================================================================
\section{Results}
\label{sec:results}

\subsection{Overall}

Over the five normal days analysed so far (8, 9, 10, 12 and 13 November), 160,046 official announcements were paired
with a bus and with the method's prediction for the same bus at the same moment: 2,854 passages of 879 trips of 72
lines at the 60 stops. Table~\ref{tab:horizon} gives the errors by how far away the bus was.

\begin{table}[t]
\centering
\caption{Errors by time to the actual passage, five normal days. \emph{Median}: median absolute error (s).
\emph{$\pm$2 min}: share of predictions within two minutes (\%). \emph{Early}: share of buses that arrive more than
two minutes \emph{before} the announced time (\%), the error that makes passengers miss the bus. The method is rounded
to whole minutes, as shown to passengers.}
\label{tab:horizon}
\small
\setlength{\tabcolsep}{4.5pt}
\begin{tabular}{lrrrrrrr}
\toprule
 & & \multicolumn{3}{c}{Official} & \multicolumn{3}{c}{Method} \\
\cmidrule(lr){3-5}\cmidrule(lr){6-8}
Minutes to arrival & Pairs & Median & $\pm$2 min & Early & Median & $\pm$2 min & Early \\
\midrule
\csname @@input\endcsname tablas/horizonte.tex 
\bottomrule
\end{tabular}
\end{table}

When the bus is less than three minutes away, both are accurate and close: 86\% of official announcements and 88\%
of the method's predictions are within two minutes. From three minutes on, they separate. With 6 to 10 minutes to go,
37\% of official announcements are within two minutes against 71\% for the method, and the median error is 167 s
against 68 s. With 10 to 20 minutes to go, the figures are 27\% and 61\%. The difference in the share within two
minutes, method minus official, has 95\% bootstrap intervals (by trip) of $[+1.2, +4.0]$ points at 0--3 minutes,
$[+22.7, +26.9]$ at 3--6, $[+32.0, +36.9]$ at 6--10, $[+31.8, +36.3]$ at 10--20 and $[+26.3, +31.1]$ at 20--30
minutes. Comparing the method in seconds instead of whole minutes changes these figures by at most one point.

\subsection{Direction of the errors}

The two sources fail in opposite directions. The official announcements are late: the median signed error is $+124$~s
at 10--20 minutes and $+149$~s at 20--30 minutes, and in about half of the announcements made 10 to 30 minutes ahead
the bus arrives more than two minutes \emph{earlier} than announced (51--52\%, against 16\% for the method). This is
the costly error for a passenger, who plans to reach the stop shortly before the announced time and finds that the bus
has gone. The method leans the other way at long horizons (median $-24$~s at 10--20 minutes and $-60$~s at 20--30
minutes): when it errs, the passenger waits a little longer at the stop rather than missing the bus.

The official announcements are not simply the published timetable: only 12--17\% of them are within one minute of the
timetable time, and the timetable itself is within two minutes of the actual passage in only 14--18\% of these cases.

\subsection{Days, stops and robustness}

The pattern holds on every day (Table~\ref{tab:days-results}): the share within two minutes is 33--38\% for the official
service and 59--72\% for the method. It holds on weekdays (37\% against 63\%) and on Sunday (33\% against 72\%), and at
the priority stops (35\% against 62\%) as at the others (37\% against 66\%). For the second announced bus ($e_2$), which
is further away, the official service is within two minutes in 26\% of cases and the method in 53\%.

\begin{table}[t]
\centering
\caption{Errors by day, all horizons up to 30 minutes. Columns as in Table~\ref{tab:horizon}. Wednesday 11 (general
strike) is shown for reference and is not part of the results of the week.}
\label{tab:days-results}
\small
\setlength{\tabcolsep}{4.5pt}
\begin{tabular}{lrrrrrrr}
\toprule
 & & \multicolumn{3}{c}{Official} & \multicolumn{3}{c}{Method} \\
\cmidrule(lr){3-5}\cmidrule(lr){6-8}
Day & Pairs & Median & $\pm$2 min & Early & Median & $\pm$2 min & Early \\
\midrule
\csname @@input\endcsname tablas/dias.tex 
\bottomrule
\end{tabular}
\end{table}

Removing the pairs in which another bus of the same line passed within two minutes of the matched one, where the
matching could pick the wrong bus, removes 1.3\% of the pairs and changes no figure by more than one point.

\subsection{The strike day}

On Wednesday 11, with minimum services, the official service was within two minutes in 36\% of cases and the method in
63\%, the same as on a normal day. The method was, however, more often late on that day: in 27\% of cases the bus
arrived more than two minutes before its prediction, against 12--16\% on normal days. Its learned travel times come from
normal traffic, and with fewer buses and lighter traffic the buses ran faster than learned.

% =====================================================================================================================
\section{Discussion}
\label{sec:discussion}

\subsection{What the comparison shows}

On the buses and stops compared, a method that uses only the operator's open data predicts arrivals more accurately
than the operator's own service from three minutes before arrival onwards, and the gap is largest, about 34 points in
the share within two minutes, when the bus is 6 to 20 minutes away. That is the range in which a passenger decides when
to leave for the stop. Below three minutes both are accurate, because both effectively see the bus approaching.

The comparison does not reveal how the official service computes its waiting times, and we do not claim to know. Its
errors grow quickly with the horizon and are biased towards announcing the bus later than it comes. That is the
behaviour one would expect from a prediction that does not learn how fast buses actually travel on each segment, which is
the core of the method \parencite{garrido2026metodo}.

\subsection{Third-party apps}

Many passengers do not use the official app but general-purpose journey planners, which receive real-time data from
operators. If such an app shows the operator's prediction unchanged, its accuracy is the official accuracy measured
here. We have not measured any third-party app directly, and their processing may differ.

\subsection{Limitations}

\begin{itemize}
  \item \textbf{Which buses are compared.} An announcement can only be compared if the actual passage of the matched
        bus can be reconstructed from the feed. This excludes buses for which the stop is at or near the end of their
        trip (the feed stops reporting them when the trip ends) and buses not reported by the feed. Of the official
        announcements of the next bus up to 25 minutes ahead, 19\% could be paired. The results describe the buses
        whose passage is observable, which are the majority of buses \emph{in the middle} of their trip, not every
        announcement.
  \item \textbf{One week, one network.} The data cover five normal days (one of them a Sunday) in November, at 60
        stops chosen for their traffic and their difficulty. Other seasons, timetables or networks may differ.
  \item \textbf{Resolution of the actual passage.} The feed changes every 60 seconds, so the actual passage is known to
        about $\pm30$~s. This adds the same noise to both sources and cannot explain differences of minutes.
  \item \textbf{Matching.} The official service does not identify the vehicle it announces. Matching by order of
        passage is what a passenger experiences, and the results do not change when ambiguous cases are removed.
  \item \textbf{The method is the replay, not the app.} The method is recomputed from the feed with its reference
        implementation, with tables frozen before the week. The live app recalibrates every night and also remembers
        buses that have recently passed, so it may perform slightly differently.
\end{itemize}

% =====================================================================================================================
\section{Reproducibility}
\label{sec:repro}

The capture program, the list of stops with the reason for each, the days used (\texttt{dias.json}), the comparison
(\texttt{comparar.py}) and the tables (\texttt{resultados.py}) are in the folder \texttt{articulo-comparacion} of the
repository. The raw captures (feed snapshots, official responses and the timetable of each day) are published as assets
of the release \texttt{comparacion-2026-11}. The method's reference implementation and the learned tables used are those
of the companion paper.

% =====================================================================================================================
\section{Conclusions}
\label{sec:conclusions}

For one week, at 60 stops of the Bizkaibus network, we recorded the official real-time waiting times minute by minute
and compared them, bus by bus, with the arrivals predicted by an open method that uses only the operator's published
positions and timetable. With the bus 6 to 20 minutes away, the official service is within two minutes of the actual
passage in about a third of cases and the method in about two thirds, and the official errors are mostly in the
direction that makes passengers miss the bus. Operators that publish their vehicle positions already make possible
predictions better than their own; adopting a learned, segment-based method, or opening the predictions to independent
evaluation, would directly benefit their passengers.
